El estudio de la evolución estelar constituye uno de los pilares fundamentales de la astrofísica moderna. Tradicionalmente, los modelos de evolución de estrellas aisladas han sido un marco robusto para comprender el ciclo de vida de los astros. Sin embargo, la vasta población de sistemas binarios y múltiples en el universo introduce escenarios complejos que desafían las aproximaciones más idealizadas. Entre estos fenómenos, las blue stragglers representan uno de los mayores enigmas: estrellas que, pese a la posición del punto de desvío de su cúmulo estelar, parecen haber encontrado una ``fuente de la eterna juventud''.

Este Trabajo de Fin de Máster nace con el propósito de arrojar luz sobre los mecanismos internos que gobiernan a estas estrellas mediante el uso de la astrosismología, una muy potente herramienta capaz de indagar en interior estelar a través de sus modos de oscilación. Tomando como base de simulación los códigos MESA y GYRE, esta investigación se sumerge en el modelado bidimensional de la estructura y pulsación de estos productos de la evolución binaria.

El alcance de este trabajo no se limita, sin embargo, a la exploración teórica abstracta. La pertinencia y el impacto de este análisis se enmarcan y justifican dentro del contexto científico de la propuesta de misión espacial HAYDN (High-precision AsteroseismologY in DeNse stellar fields) de la Agencia Espacial Europea (ESA). Diseñada para operar en campos estelares densos y cúmulos jóvenes mediante fotometría de altísima precisión, HAYDN promete revolucionar nuestra comprensión de la física estelar. Los modelos teóricos desarrollados en esta tesis pretenden servir como un primer paso de contribución al marco preparatorio de dicha misión, demostrando el potencial que la astrosismología tiene para revelar la historia evolutiva y estructural de este tipo de estrellas.

Afrontar este proyecto ha supuesto no solo un reto intelectual en el ámbito de la astrofísica estelar, sino también un profundo aprendizaje en el desarrollo de herramientas computacionales, la modificación de entornos de simulación numérica y la interpretación de espectros de frecuencia pulsacionales. Es el resultado de meses de dedicación orientados a un horizonte donde la teoría astrosismológica y el futuro de la exploración espacial europea se dan la mano.